\documentclass[12pt]{article}
\usepackage[utf8]{inputenc}
\usepackage[T1]{fontenc}
\usepackage{hyperref}
\title{Samre Language Resources for K-12 Curriculum Development}
\author{Compiled by Claude Sonnet 5.5 (claude-sonnet-5-5) for Andr\'as Kornai}
\date{}
\begin{document}
\maketitle

\section{Introduction}
Samre (ISO 639-3 sxm) is a dormant Pearic (Austro-Asiatic) language of Cambodia and Thailand. Search results describe it as nearly extinct, with a 1998 survey reporting 20--30 speakers in Nonsi Subdistrict, Bo Rai District, Trat Province, Thailand. Sources disagree on naming and location (Samre, Somray, Samre of Pursat), so care is needed not to mix up varieties. Very little material exists, and none found is a school curriculum.

\section{Primary Educational Resources}
No primers, textbooks or school materials for Samre were found in the searches.

\section{Linguistic and Reference Materials}
\begin{itemize}
\item Wikipedia, ``Samre language of Pursat'': \url{https://en.wikipedia.org/wiki/Samre_language_of_Pursat}
\item Wikipedia, ``Pearic languages'' (family overview): \url{https://en.wikipedia.org/wiki/Pearic_languages}
\item Headley, Robert K. (1977), ``A Pearic Vocabulary'', \emph{Mon-Khmer Studies} 6; a word list covering Samre varieties among other Pearic languages (bibliographic record on Glottolog; contents not read): \url{https://glottolog.org/resource/reference/id/103665}
\item Thongkum, Theraphan (1984), ``The Samre language'' (Samre of Pursat; bibliographic record only): \url{https://glottolog.org/resource/reference/id/306421}
\item Glottolog language record for the related Somray entry (iso smu): \url{https://glottolog.org/resource/languoid/iso/smu}
\end{itemize}

\section{Religious and Community Resources}
No religious or community-organization resources specific to Samre were found.

\section{Digital and Audio Resources}
\begin{itemize}
\item Endangered Languages Project entry (a catalogue record; not verified to contain audio): \url{https://www.endangeredlanguages.com/lang/1223}
\item Lexvo entry for code sxm (metadata only): \url{https://finto.fi/lexvo/en/page/sxm}
\end{itemize}
No audio or video recordings were found.

\section{Curriculum Development Context}
The language is dormant and not known to be taught. A revitalization curriculum would have to be reconstructed from archival word lists and linguistic descriptions such as those above, ideally with the descendant community in Thailand. Only a vocabulary-level, heritage-awareness programme appears feasible on current evidence. The resources above were not examined in full, and their coverage should be checked before use.

\end{document}